\chapter{Deciding Bitvector Formulas by Bitblasting}
\label{ch:background-bv}
\epigraph{The computer is not a cheat, but an instrument of exploration.}{Kenneth Appel and Wolfgang Haken, on the Four Color Theorem}

\paragraph*{Attribution.} Parts of this chapter, in particular the account of
Lean's bitvector library and the \bvdecide{} tactic, are based on
\emph{Interactive Bitvector Reasoning using Verified
Bit-Blasting}~\citep{leanbv} (OOPSLA 2025), joint work led by Henrik Böving on
which I am a second author. The decision procedures and the
floating-point bitblaster developed in the chapters that follow are all built
on top of the foundations of fixed-width bitblasting (\cref{ch:mono-width,ch:multi-width,ch:floating-point}).

\lettrine{B}{itvectors} are the lingua franca of low-level verification.
The pointer-free fragment of a programming language can be modelled by
bitvectors, matching the memory representation of integers, enums, and
structures. Bitvectors also model integer operations in compiler IRs such
as LLVM-IR~\citep{zhao2012formalizing}, machine
code~\citep{armstrong2019isa,reid2016trustworthy}, combinational logic in
hardware~\citep{bourgeat2020essence}, and are the substrate on which
floating-point arithmetic is defined~\citep{boldo2011flocq}.
Here, we review how SMT solving over the theory of fixed-width bitvectors (\qfbv{})
is performed by reducing the problem to a SAT problem,
followed by a verified certificate check.
This work is integrated into Lean's verified bitblaster, \bvdecide{}.
\bvdecide{}.

\section{SMT Solving over Bitvectors}
\label{sec:bg:smt}

Satisfiability Modulo Theories (SMT) can be viewed as a generalisation of Boolean satisfiability (SAT).
In Boolean satisfiability, one asks whether a propositional formula is satisfiable.
That is, we have Boolean variables $x_1, \dots, x_n$ and a formula $\phi(x_1, \dots, x_n)$ built from them using the logical connectives $\land$, $\lor$, and $\lnot$,
and we ask whether there exists an assignment of truth values to the variables such that $\phi$ evaluates to true.
SMT generalises this by allowing variables to range over richer domains,
such as integers, real numbers, or bitvectors, and allowing formulae to include more operations, such as arithmetic or bitwise operations.
The data that corresponds to the richer domains, and the operations that can be performed on them, is called a theory.
A satisfying assignment for an SMT formula is an assignment of values to the variables such that the formula evaluates to true under the interpretation of the theory.
This is called a model of the formula. We will not go into the details of SMT solving here,
but refer the reader to \citet{kroening2016decision} for the algorithmics of SMT solving,
to \citet{hodges1997shorter} for model theory, and \citet{boolos2002computability} for my favourite introduction to logic.
\emph{Computability and Logic} also has an extremely slick proof of G\"odel's incompleteness as a corollary
of the halting problem, which I found very illuminating.

SMT solvers such as Z3~\citep{de2008z3},
cvc5~\citep{barrett2011cvc4,barbosa2022cvc5}, and Bitwuzla~\citep{niemetz2023bitwuzla}
decide the satisfiability of these first-order formulae. The SMT-LIB standard~\citep{barrett2010smt,smtlib}
specifies these theories. The theory relevant to this chapter is \qfbv{}, the
\emph{quantifier-free theory of fixed-width bitvectors}. \qfbv{} exposes a
large vocabulary of operations over bitvectors of a fixed, concrete width (say, $64$):
modular arithmetic (addition, multiplication, signed and unsigned division and
remainder), bitwise Boolean operations, shifts, comparisons, concatenation, extraction, and overflow predicates.
Every \qfbv{} formula fixes the width of each of its variables to a concrete numeral.

Because a bitvector of width $w$ ranges over the finite domain $\{0,1\}^w$,
\qfbv{} admits an effective reduction to Boolean satisfiability.
Here, each bitvector variable is replaced by $w$ Boolean variables,
and each operation by the Boolean circuit that computes it.
For example, one builds a ripple-carry adder for addition,
or a shift-and-add circuit for multiplication.
To prove a universally quantified fixed-width statement, such as $\forall x, y : \bv{64}, x + y = y + x$,
one first negates it to get an existentially quantified statement $\exists x, y : \bv{64}, x + y \neq y + x$.
Intuitively, we are looking for a counterexample to the original statement, and if we find one, we have disproved it.
If we can find no counterexample, then the original statement is a theorem.

To find counterexamples, we convert the bitvector problem into a Boolean satisfiability problem by the bitblasting approach mentioned above.
One then asks a SAT solver whether the resulting SAT formula/circuit is satisfiable.
If the formula is satisfiable, we have found a counterexample!
On the other hand, if the formula is \emph{unsatisfiable}, then the original statement is a theorem.
Now, this immediately raises the question of trust: when the solver claims a formula is unsatisfiable, how do we know it is correct?
After all, for satisfiability, it is trivial, since we can plug the claimed solution back into the original formula and check that it evaluates to true.
However, this doesn't work for unsatisfiability, since we would have to check all possible assignments to the variables, which is infeasible.
Intuitively, this is because while $\mathrm{SAT}$ is in $\mathrm{NP}$, which is defined by having a polynomially verifiable certificate,
$\mathrm{UNSAT}$ is in $\mathrm{coNP}$, which is defined by having a polynomially verifiable refutation.
This brings us to the wonderful world of UNSAT certificates and proof complexity,
which are ``compact proof traces'' that the formula is unsatisfiable.

\paragraph*{UNSAT Certificates.}

Modern SAT solvers emit \emph{UNSAT certificates},
which convince a checker that the input formula is unsatisfiable by deriving the empty clause (a contradiction) from the input CNF.
The de facto standards here are the DRAT~\citep{heule2016drat} and LRAT~\citep{cruz2017efficient} formats.
The DRAT format is a sequence of derivation steps
in which each step records only the clause derived,
leaving the checker to rediscover why that step is sound.
state-of-the-art solver \cadical{}~\citep{biere2024cadical} can emit LRAT
certificates directly.
Checking certificates removes the SAT solver from the trusted code base.
A buggy solver can, at worst, produce a certificate that the checker rejects.

A lot of research has gone into designing excellent UNSAT certificate formats for SAT solvers.
LRAT is expressive enough to express all the preprocessing tricks that SAT solvers perform,
while being concise enough to be checked efficiently.
Unfortunately, no comparable mature certificate format exists at the \emph{SMT} level for bitvectors~\citep{barrett2015proofs},
which is why we perform bitblasting inside the proof assistant and outsource only the SAT search to an unverified solver.

% \paragraph*{And-Inverter Graphs.} The circuits produced by bitblasting are
% conventionally represented as \emph{And-Inverter Graphs} (AIGs): directed
% acyclic graphs whose nodes are inputs, constants, and binary AND gates, with
% optional negation markers on gate inputs. Despite this minimal gate set, AIGs
% can represent every Boolean circuit, and their simplicity makes them amenable
% to \emph{structural hashing} (sharing identical subcircuits) and cheap local
% rewrites~\citep{brummayer2006local}. SAT solvers, in turn, accept formulas in
% conjunctive normal form (CNF), so the AIG is translated to CNF by the Tseitin
% transformation~\citep{tseitin1983complexity}, which introduces one fresh
% variable per gate and produces a CNF that is \emph{equisatisfiable} with the
% circuit, at only linear size overhead.


\section{Lean's Bitvector Library}
\label{sec:bg:bitvec}

Lean's standard library provides a canonical bitvector type,
\texttt{BitVec w}, whose design was driven by the needs of proof automation
and, in particular, by the proofs required to build a verified
bitblaster~\citep{leanbv} (\cref{sec:bg:bvdecide}). The library covers all
operations of SMT-LIB 2.7 \qfbv{}, including division, remainder, and the
overflow predicates, and comprises roughly 150 definitions and over 900
theorems. We recall the design decisions that later chapters rely on.

\subsection{Definition and Constructors}

The low-level definition of \texttt{BitVec w} is a natural number $n$ together
with a proof $\texttt{isLt} : n < 2^w$. This representation, rather than the alternative of a vector of Booleans,
permits an efficient implementation of all operations on the bitvector.
Recall that the Lean kernel uses GMP for fast \texttt{Nat} arithmetic,
so the bitvector library uses this to get access to fast operations on bitvectors.


The representation is \emph{signless}, so a bitvector carries no notion of being signed or unsigned.
Instead, individual operations on the bitvector interpret it as one or the other.
This matches the semantics of SMT-LIB and of compiler IRs such as LLVM,
and also permits some fun hacker koans, where we interpret a signed number as an unsigned value for some cute effect.
For example, one can perform a two-sided range check of $0 \leq x \leq 128$ (say $x$ is a 32-bit signed integer)
by checking that $x$ is \emph{unsigned} less than or equal to $128$.
If $x$ is a negative number, then it will be interpreted as a large unsigned number, and the check will fail!

To go from natural numbers to bitvectors and vice versa,
any natural number coerces to a bitvector by reduction modulo $2^w$:
$\texttt{ofNat}\ (n : \N) \mapsto (n \bmod 2^w) : \BV{w}$. Conversely,
\texttt{x.toNat} recovers the underlying natural number.

\subsection{Bitwise Operations}

Bitwise operations reuse Lean's bitwise operations on natural numbers. The
value of an individual bit is retrieved with \texttt{v.getLsbD i} (``get the
\texttt{i}-th least significant bit, with default value \texttt{0}'').
Crucially, \texttt{getLsbD} is \emph{total}: it returns \texttt{false} for out-of-range
indices instead of demanding a bounds proof, which avoids the repeated
in-bounds obligations a list-based definition would incur.

\subsection{Arithmetic Operations}

Addition, subtraction, and multiplication are arithmetic modulo $2^w$; the
result is less than $2^w$ by construction, discharging \texttt{isLt}.
Sign-sensitive operations (signed division, sign extension, arithmetic shift right)
are defined via coercions to and from \Z{}:
\texttt{BitVec.toInt} computes the two's-complement value $2^0 b_0 + \dots + 2^{w-2}b_{w-2} - 2^{w-1}b_{w-1}$, and \texttt{BitVec.ofInt}
maps an integer $i$ to $\texttt{BitVec.ofNat}\ (i \bmod 2^w)$ using the
Euclidean remainder, which lies in $[0, 2^w)$.

These closed-form definitions are convenient to reason with but bear no
resemblance to the circuits a bitblaster emits. The library therefore also
provides \emph{recursive models} of those circuits.
For example, \texttt{adc} for carry-propagating addition, \texttt{mulRec} for shift-and-add multiplication,
all of which are proven equal to the corresponding closed-form operation at every width.

\section{\texorpdfstring{\bvdecide{}}{bv\_decide}: Verified Bitblasting for Fixed Widths}
\label{sec:bg:bvdecide}

\bvdecide{}~\citep{leanbv} is Lean's decision procedure for the universally
quantified, fixed-width theory of bitvectors, with full support for the
SMT-LIB \qfbv{} operations. Rather than trusting an external SMT solver,
which would add the solver's preprocessing, bitblasting, and (typically uncertified) SAT reasoning to the trusted code base,
\bvdecide{} implements a \emph{verified} bitblaster inside Lean and outsources only the SAT search,
whose UNSAT answers are re-checked against LRAT certificates (\cref{sec:bg:smt}).
Because realistic UNSAT proofs are far too large for kernel type checking,
the bitblaster and certificate checker do not produce proof terms; instead they run as compiled code via proof by reflection.
The user thus trusts the Lean compiler in addition to the kernel, which is a much smaller leap of faith than trusting a full SMT solver.
Also, a moment's thought shows that one already trusts the Lean compiler,
by using Lean as a programming language, and so it is not an extreme assumption to make that the compiler is correct.

To prove a fixed-width goal, \bvdecide{} performs the following steps
(\cref{fig:bg:bvdecide}):
\begin{enumerate}
  \item Initiate a proof by contradiction, replacing the goal with
  \texttt{False} and adding the goal's negation as a hypothesis.
  \item Simplify and canonicalize all fixed-width bitvector hypotheses
  (the \bvnormalize{} pass).
  \item Bitblast the conjunction of the simplified hypotheses into an
  equisatisfiable AIG, and translate the AIG to CNF.
  \item Hand the CNF to the SAT solver \cadical{}. A SAT answer yields a
  counterexample presented to the user; an UNSAT answer yields an LRAT
  certificate, which a verified checker validates inside Lean.
  \item Since every step is proven to preserve satisfiability, the tactic
  concludes that the hypotheses are inconsistent, finishing the proof.
\end{enumerate}

% \subsection{Intrinsically Well-Typed AIGs}
% \label{sec:bg:bvdecide-aig}
% 
% \begin{figure}
% \begin{leanfootnotesize}
% inductive Decl (α : Type) where
%   | const (b : Bool)
%   | atom (idx : α)
%   | gate (l r : Nat) (linv rinv : Bool)
% 
% def IsDAG (α : Type) (ds : Array (Decl α)) :=
%   ∀ {i lhs rhs linv rinv} (h : i < ds.size),
%     ds[i] = .gate lhs rhs linv rinv → lhs < i ∧ rhs < i
% 
% structure AIG (α : Type) [DecidableEq α] [Hashable α] where
%   decls : Array (AIG.Decl α)     -- the circuit
%   cache : AIG.Cache α decls      -- verified cache for subterm sharing
%   invariant : AIG.IsDAG α decls  -- acyclicity
% \end{leanfootnotesize}
% \caption{Core AIG data structures (simplified). Acyclicity is encoded by the
% proposition \texttt{AIG.IsDAG}, which forces gates to point only to
% lower-indexed nodes; the cache is a verified hash map that is well-formed with
% respect to the declaration array, enabling full subterm sharing while
% remaining correct by construction. Adapted from the \bvdecide{}
% paper~\citep{leanbv}.}
% \label{fig:bg:aig-data-structures}
% \end{figure}
% 
% The bridge between bitvector expressions and the SAT solver is an AIG
% (\cref{sec:bg:smt}). Lean, being purely functional, cannot build DAGs in the
% traditional pointers-to-heap-values style. Instead, an AIG is an \emph{array}
% of declarations---constants, input atoms, and AND gates whose inputs are
% array indices, each with an optional negation flag
% (\cref{fig:bg:aig-data-structures}). Acyclicity is captured intrinsically by
% the invariant that a gate's inputs have strictly smaller indices than the gate
% itself. The AIG additionally carries a verified hash map from declarations to
% indices, enabling structural sharing; its well-formedness invariant states
% that it only maps declarations to the index at which they actually occur. The
% construction is heavily inspired by AIGNET~\citep{davis2013verified}. Thanks
% to Lean's functional-but-in-place (FBIP) reuse of uniquely-referenced values,
% keeping a single live copy of the AIG means the
% array and hash map are mutated in place, so the purely functional code has
% imperative performance. The gate constructor is a \emph{smart constructor}
% that applies verified versions of the level-1 local
% rewrites~\citep{brummayer2006local} and orders gate inputs by index to improve
% sharing; it is proven to produce results equisatisfiable with the naive
% constructor.
% 
% Operations built on top of these primitives---complex gates, arithmetic
% circuits---are verified through the notion of a \emph{lawful operator}: a
% function that takes an AIG (plus references into it) and returns an extended
% AIG together with a reference to its output, without shrinking the array or
% modifying existing entries. Lawfulness enables local reasoning: when reasoning
% about a subcircuit, all later lawful extensions can simply be forgotten, and
% since compositions of lawful operators are lawful, lawfulness is proven once
% for the primitives and inherited by everything built from them.

\subsection{Verified Bitblasting}
\label{sec:bg:bvdecide-bitblasting}

\begin{figure}
	\includegraphics[width=\textwidth]{chapters/background-bv/images/bitvec-paper.pdf}
	\caption{\bvdecide{} provides an end-to-end verified bitblasting
  pipeline inside Lean. The trusted code base for the respective steps is
  listed in red; when rewriting alone proves a theorem, there is no need to
  trust the compiler. Figure reproduced from the \bvdecide{}
  paper~\citep{leanbv}.}
	\label{fig:bg:bvdecide}
\end{figure}

After simplification, all hypotheses consisting only of fixed-width bitvector
or Boolean statements are collected; subterms with custom definitions are
abstracted as uninterpreted bitvector atoms, which preserves at least basic
reasoning for user-defined functions. The conjunction of the hypotheses is
then bitblasted into an equisatisfiable AIG using circuit constructions from
Bitwuzla and Z3~\citep{de2008z3}. Bitwise operations bitblast to trivial
gate-per-bit circuits verified directly from the AIG's evaluation semantics.
Arithmetic operations require genuine circuits. First, the model (such as \texttt{adc} or \texttt{mulRec})
is proven equal to the operation at all widths.
Then, by induction, each recursive step of the bitblaster is shown to add a subcircuit to the AIG that mirrors the
corresponding step of the model. For performance, the bitblaster also maintains a verified
expression-level cache mapping bitvector expressions to vectors of AIG
references, using Lean's pointer-equality primitives~\citep{Selsam2020} on
maximally shared expression nodes for fast lookups.

\subsection{Certified UNSAT Answers via LRAT Checking}
\label{sec:bg:bvdecide-sat}

\begin{figure}
\begin{leanfootnotesize}
-- Parse and check an LRAT proof.
def verifyCert (cnf : CNF Nat) (cert : LratCert) : Bool :=
  match LRAT.parseLRATProof cert with
  | .ok lratProof => LRAT.check lratProof cnf
  | .error _ => false

-- `verifyCert` returning true implies that the CNF is UNSAT.
theorem verifyCert_correct : ∀ cnf cert, verifyCert cnf cert = true → cnf.Unsat

-- Verify a BV expression is UNSAT by using the certificate to witness UNSAT.
def verifyBVExpr (bv : BVLogicalExpr) (cert : LratCert) : Bool :=
  verifyCert (AIG.toCNF bv.bitblast.relabelNat) cert

-- The computable verifier producing `true` implies that the formula is UNSAT.
theorem unsat_of_verifyBVExpr_eq_true (bv : BVLogicalExpr) (c : LratCert)
    (h : verifyBVExpr bv c = true) : bv.Unsat
\end{leanfootnotesize}
\caption{The verified LRAT checker and its connection to the bitblaster. This
enables Boolean reflection: the computable hypothesis
\texttt{verifyBVExpr bv c = true}, whose proof is \texttt{Eq.refl}, implies
the proposition \texttt{bv.Unsat}. Adapted from the \bvdecide{}
paper~\citep{leanbv}.}
\label{fig:bg:unsat-cert-check}
\end{figure}

The AIG is translated to CNF by a Tseitin transformation proven to produce an
equisatisfiable formula, and the CNF is passed to
\cadical{}~\citep{biere2024cadical}. If \cadical{} finds a satisfying
assignment, \bvdecide{} maps it back to an assignment of the original
bitvector variables and reports it as a counterexample. If \cadical{} reports
UNSAT, it emits an LRAT certificate~\citep{cruz2017efficient}, which is parsed
(by a port of the binary parser of \texttt{lrat-trim}) and trimmed, typically
dropping about half of the steps. Neither parsing nor trimming needs to be
verified: they merely prepare the input for the \emph{verified} LRAT checker,
and a bug in them can only cause the checker to fail, never to accept a wrong
proof. Running the checker establishes that the CNF, and hence, by the chain
of equisatisfiability proofs, the original conjunction of hypotheses, is
UNSAT (\cref{fig:bg:unsat-cert-check}), completing the contradiction proof.

\section{Discussion}

Overall, \bvdecide{} provides a fully verified decision procedure for the quantifier-free theory of fixed-width bitvectors.
This gives us the ability to automate proving fixed-width bitvector formulas in Lean.
This is already very powerful, and lets us automate away the correctness of peephole rewrites at a fixed width.
However, we'd like to go further, and automate reasoning about bitvector formulas at an \emph{arbitrary} width,
as well as reasoning about floating-point formulas, which are built on top of bitvectors.
We will explore these extensions in the following chapters.

